\documentclass[11pt]{article}
% Shared preamble for Vietnamese companion manuscripts
\usepackage{fontspec}
\setmainfont{DejaVu Serif}
\usepackage[a4paper,margin=2.1cm]{geometry}
\usepackage{microtype}
\usepackage{amsmath,amssymb}
\usepackage{booktabs,tabularx,array,longtable}
\usepackage{graphicx}
\usepackage{enumitem}
\usepackage[hidelinks]{hyperref}
\usepackage{xurl}
\usepackage{titlesec}
\usepackage{fancyhdr}
\setlength{\parindent}{0pt}
\setlength{\parskip}{0.55em}
\setlist[itemize]{leftmargin=1.5em,itemsep=0.25em}
\setlist[enumerate]{leftmargin=1.7em,itemsep=0.25em}
\renewcommand{\arraystretch}{1.12}
\titleformat{\section}{\large\bfseries}{\thesection}{0.6em}{}
\titleformat{\subsection}{\normalsize\bfseries}{\thesubsection}{0.5em}{}
\pagestyle{fancy}\fancyhf{}\fancyfoot[C]{\small Bản tiếng Việt đồng hành -- trang \thepage}\renewcommand{\headrulewidth}{0pt}

\emergencystretch=3em
\sloppy

\title{\textbf{Từ bối cảnh sức khỏe dọc đến hành động an toàn: hợp đồng quản trị đồng phiên bản với rào chắn an toàn Tầng~1 phi-LLM cho AI y học gia đình}}
\author{Nguyễn Ngọc Thiện, Trịnh Minh Quang\\
Trường THPT Hai Bà Trưng, Huế, Việt Nam\\
Tác giả liên hệ: \texttt{kakangocthien109@gmail.com}}
\date{Tháng 8 năm 2026}
\begin{document}
\maketitle
\thispagestyle{empty}

\textbf{Chuyên đề:} AI và chuyển đổi số trong Y tế và Giáo dục Y khoa

\section{Đặt vấn đề trong Y học Gia đình và Chăm sóc Ban đầu}
Trong y học gia đình và chăm sóc sức khỏe ban đầu, hồ sơ bệnh án điện tử tích lũy theo chiều dọc qua nhiều năm và bao quát đa dạng các đợt khám: quản lý bệnh mạn tính (đái tháo đường, tăng huyết áp), tiền sử dùng thuốc phức tạp (polypharmacy), lịch tiêm chủng, sàng lọc định kỳ, cùng các thông tin tâm lý -- xã hội nhạy cảm. Khi ứng dụng trí tuệ nhân tạo (AI) và mô hình ngôn ngữ lớn (LLM) hỗ trợ bác sĩ gia đình, các hệ thống có trạng thái đối mặt với hai thách thức an toàn cốt tử:

\begin{enumerate}[leftmargin=1.8em]
\item \textbf{Lệch pha thời gian từ đọc đến ghi (TOCTOU Drift):} Mô hình có thể đọc một phần hồ sơ bệnh nhân khi dữ liệu và quyền truy cập còn hợp lệ tại thời điểm $t_1$, nhưng chỉ trả về đề xuất hành động lâm sàng (kê đơn, lên lịch tái khám, cập nhật hồ sơ) tại thời điểm $t_3$, sau khi hồ sơ, trạng thái đồng thuận của bệnh nhân hoặc chính sách lâm sàng đã thay đổi tại $t_2$. Một lần đọc hợp lệ vì thế không còn đủ điều kiện để cho phép một thao tác ghi dữ liệu về sau.
\item \textbf{Nguy cơ lộ lọt thông tin nhạy cảm (Over-disclosure) và ảo giác số học thời gian:} Việc nạp toàn bộ bệnh án không giới hạn vào LLM vi phạm nghiêm trọng nguyên tắc giảm thiểu dữ liệu (\emph{Zero-PHI over-disclosure}), làm lộ các dữ liệu nhạy cảm không liên quan (sức khỏe sinh sản, tâm thần, yếu tố xã hội). Đồng thời, LLM thường xuyên mắc ảo giác khi tính toán thời hạn khám sàng lọc, khoảng cách liều và phát hiện tương tác thuốc (DDI).
\end{enumerate}

\section{Kiến trúc Rào chắn Trạng thái Hai Tầng (Dual-Layer State Barrier) và Hợp đồng GLHS}
Để giải quyết triệt để các rủi ro trên, chúng tôi phát triển kiến trúc GLHS (Governed Longitudinal Health Snapshot) trong hệ thống CLARA-Care với cơ chế bảo vệ kép:

\begin{itemize}[leftmargin=1.8em]
\item \textbf{Rào chắn an toàn Tầng~1 phi-LLM (Deterministic Reconciliation Kernel tại API Gateway):} Một động cơ tất định chạy bằng Python/PostgreSQL thực thi toàn bộ các ràng buộc lâm sàng cứng trước khi chuyển dữ liệu cho AI:
\begin{itemize}
  \item Thực hiện phép toán thời gian hai chiều nghiêm ngặt ($\tau_{\text{valid}} \cap \tau_{\text{known}}$) để ngăn ngừa nhiễm bẩn dữ liệu hồi cứu.
  \item Tính toán cửa sổ thời hạn tái khám, lịch tiêm chủng và thời gian ân hạn (grace period) theo đúng phác đồ y học gia đình bằng phép toán tất định (loại bỏ hoàn toàn ảo giác tính toán số học thời gian của LLM).
  \item Đối soát tự động ma trận tương tác thuốc nguy cơ cao (DDI) và chống chỉ định điều trị dựa trên các bộ mã chuẩn (RxNorm, SNOMED CT, LOINC).
  \item Xác thực trạng thái đồng thuận động của bệnh nhân, phân quyền RBAC và phiên bản chính sách quản trị.
  \item Đóng gói và tạo mã băm Merkle ($\Pi_{\text{Merkle}}$) thành Bản chụp trạng thái sức khỏe giới hạn theo tác vụ (\textbf{Task-Bounded Health State Snapshot -- THSS}), thực thi nguyên tắc giảm thiểu dữ liệu theo tác vụ (chỉ trích xuất đúng lát cắt dữ liệu cần thiết cho tác vụ lâm sàng cụ thể).
\end{itemize}
\item \textbf{Tầng~2 (Epistemic Arbiter -- LLM):} Mô hình ngôn ngữ lớn chỉ tập trung suy luận ngữ nghĩa định tính, tổng hợp bệnh sử và giải thích lâm sàng, hoàn toàn được cách ly khỏi các phép tính số học thời gian thô và quy tắc an toàn cứng.
\item \textbf{Chuyển trạng thái có quản trị (Governed State Transition -- GST):} Đề xuất cập nhật do LLM tạo ra bắt buộc phải mang theo mã định danh THSS đã sử dụng. Trước khi ghi dữ liệu vào cơ sở dữ liệu bệnh án, GST tự động kiểm tra lại độ mới của trạng thái và thẩm quyền quản trị. Nếu dữ liệu hoặc đồng thuận đã bị thay đổi trong lúc mô hình suy luận, giao dịch ghi sẽ bị từ chối ngay lập tức (fail-closed).
\end{itemize}

\section{Đánh giá Tương quan Quyết định Lâm sàng và Giảm thiểu Bối cảnh ($N=384$)}
Nhằm đánh giá thực nghiệm câu hỏi: \emph{``Việc giảm thiểu bối cảnh qua Strict THSS ảnh hưởng như thế nào đến chất lượng ra quyết định lâm sàng so với việc nạp toàn bộ bệnh án?''}, chúng tôi thực hiện phân tích đối chiếu ghép cặp cấp đối tượng trên đoàn hệ độc lập gồm 384 đối tượng y học gia đình ($N=384$, 3.456 ô thực thi).

Kết quả đánh giá đối đầu:
\begin{itemize}[leftmargin=1.8em]
\item \textbf{Kết quả ghép cặp (Sign-Test):} Trên 384 đối tượng, ghi nhận 70 ca Strict tốt hơn, 73 ca Full tốt hơn và 241 ca hòa ($\hat{\Delta} = -0{,}781\%$, $SE = 3{,}12\%$, kiểm định dấu hai phía $p = 0{,}8672$).
\item \textbf{Khoảng tin cậy Bootstrap 95\%:} Khoảng tin cậy thực nghiệm cho độ chênh lệch là $[-6{,}77\%, +5{,}21\%]$. Do khoảng tin cậy này cắt qua biên độ hẹp $\pm 2{,}0\%$, nghiên cứu ở quy mô $N=384$ chưa đủ lực thống kê để xác nhận tính không thua kém hoặc tương đương thống kê biên hẹp (đòi hỏi cỡ mẫu $N \ge 7.500$ đối tượng).
\item \textbf{Ý nghĩa thực tiễn:} Thay vì tuyên bố không suy giảm chất lượng lâm sàng tuyệt đối hay tối ưu Pareto, kết quả chỉ ra rằng độ lệch chất lượng quyết định nằm trong phạm vi phương sai thực nghiệm, trong khi việc giảm thiểu bối cảnh giúp giảm 87{,}4\% lượng token tiêu thụ (trung bình 412 so với 3.280 tokens, $p < 10^{-40}$) và giảm 68{,}2\% độ trễ suy luận.
\item \textbf{Giảm thiểu dữ liệu và an toàn trạng thái:} Hệ thống đảm bảo giảm thiểu phơi nhiễm dữ liệu ngoài phạm vi tác vụ (0{,}0\% lộ lọt thực thể nhạy cảm trong bộ kiểm thử) và ngăn chặn 256/256 lần ghi sai TOCTOU ($p = 1{,}727 \times 10^{-77}$ so với baseline không ràng buộc).
\end{itemize}

\section{Kết quả Đánh giá Hệ thống Bổ trợ}
\begin{enumerate}[leftmargin=1.8em]
\item \textbf{Đoàn hệ mô hình 64 đối tượng:} Đánh giá trên Claude Sonnet 4.6 và Gemini 3.6 Flash High ($1.152$ ô đánh giá lặp lại) cho thấy Strict THSS đạt đúng tuyệt đối trên mọi trục tiêu chí ở 63/64 đối tượng với Claude (98{,}44\%) và 64/64 với Gemini (100\%).
\item \textbf{Ma trận đồng thời PostgreSQL TOCTOU (12 lịch):} Thử nghiệm các tình huống xung đột phức tạp (thu hồi đồng thuận, hạ quyền bác sĩ, cập nhật phiên bản chính sách, sửa đổi bệnh án đồng thời) ghi nhận $0$ thao tác ghi sai trái phép; toàn bộ các đề xuất dựa trên dữ liệu cũ bị từ chối chính xác trong khi các thao tác hợp lệ được chấp thuận an toàn.
\item \textbf{Kiểm tra mô hình hữu hạn (Formal Bounded Exploration):} Duyệt 21.361 trạng thái và 90.432 bước chuyển đến độ sâu 5, xác nhận không có bất kỳ vi phạm nào trên 11 bất biến an toàn toán học.
\end{enumerate}

\begin{table}[H]\centering\small
\caption{Tổng hợp kết quả đánh giá hệ thống GLHS và rào chắn Tầng~1 cho Y học Gia đình.}
\vspace{0.4em}
\begin{tabularx}{\textwidth}{p{0.28\textwidth}p{0.24\textwidth}X}
\toprule
\textbf{Thành phần đánh giá} & \textbf{Quy mô / Phương pháp} & \textbf{Kết quả \& Ý nghĩa lâm sàng} \\
\midrule
Đối chiếu cấp đối tượng & $N=384$ đối tượng, 3.456 ô & $\hat{\Delta} = -0{,}781\%$ ($p=0{,}867$); KTC bootstrap 95\% $[-6{,}77\%, +5{,}21\%]$ \\
Giảm thiểu dữ liệu & Đo lường thực thể nhạy cảm & 0{,}0\% lộ lọt PHI ngoài phạm vi tác vụ; giảm 87{,}4\% token tiêu thụ \\
Rào chắn Tầng~1 phi-LLM & $N=1.000$ phép tính thời hạn & Tính toán số học thời gian tất định, loại bỏ ảo giác tiêm chủng/tái khám \\
Ngăn chặn xung đột TOCTOU & Matched ablation ($N=320$) & Chặn 256/256 lần ghi sai ($p=1{,}727 \times 10^{-77}$); 0 vi phạm trên 12 lịch PostgreSQL \\
Duyệt bất biến trạng thái & Formal model checking & 21.361 trạng thái, 90.432 chuyển trạng thái, 0 lỗi trên 11 bất biến \\
\bottomrule
\end{tabularx}
\end{table}

\section{Thảo luận, Ranh giới Tuyên bố và Kết luận}
\textbf{Ranh giới tuyên bố an toàn:} GLHS được thiết kế như một cơ chế tin học y tế quản trị tính nhất quán và bảo mật luồng dữ liệu đọc--ghi, \emph{không phải là một thiết bị y tế hay hệ thống chẩn đoán tự động thay thế vai trò chuyên môn của bác sĩ gia đình}. Bằng chứng hiện tại dựa trên đoàn hệ tổng hợp và đo đạc kỹ thuật phần mềm; việc triển khai thực tế đòi hỏi các thử nghiệm lâm sàng tiến cứu được hội đồng đạo đức y sinh phê duyệt.

\textbf{Kết luận:} GLHS chứng minh rằng việc kết hợp rào chắn an toàn Tầng~1 phi-LLM với hợp đồng đồng phiên bản từ công bố dữ liệu đến ghi thay đổi giúp giảm thiểu phơi nhiễm dữ liệu người bệnh và ngăn chặn rủi ro sai lệch dữ liệu trong y học gia đình, trong khi vẫn duy trì chất lượng suy luận lâm sàng trong phạm vi phương sai thực nghiệm.

\textbf{Từ khóa:} tin học y học gia đình; AI chăm sóc ban đầu; giảm thiểu dữ liệu; bối cảnh theo tác vụ; rào chắn an toàn phi-LLM; nhất quán trạng thái; quản trị lâm sàng.
\end{document}
